\documentclass[10pt]{beamer}

% Tema Limpo e Profissional
\usetheme{Madrid}
\usecolortheme{whale}

% Pacotes Essenciais
\usepackage[utf8]{inputenc}
\usepackage[portuguese]{babel}
\usepackage{graphicx}
\usepackage{booktabs}
\usepackage{amsmath}
\usepackage{amssymb}
\usepackage{listings}
\usepackage{xcolor}

% Configuracoes de Codigo (Python)
\lstset{
    language=Python,
    basicstyle=\ttfamily\tiny,
    keywordstyle=\color{blue},
    stringstyle=\color{red},
    commentstyle=\color{green!60!black},
    breaklines=true,
    showstringspaces=false
}

% Metadados
\title[IFCE Tauá - Planejamento Ágil]{Simulação de Reunião de Backlog e Estimativas}
\subtitle{Projeto FamilyFinance.AI: Finanças Familiares com Inteligência Artificial}
\author[Prof. Reginaldo Fernandes]{Prof. Reginaldo Fernandes\texorpdfstring{\\ \small{Engenharia de Software}}{}}
\institute[IFCE]{Instituto Federal do Ceará -- Campus Tauá}
\date{\today}

\begin{document}

\begin{frame}
    \titlepage
\end{frame}

\begin{frame}{Sumário da Aula}
    \tableofcontents
\end{frame}

\section{Introdução e Problema}
\begin{frame}{O Problema de Negócio: Projeto FamilyFinance.AI}
    \begin{columns}
        \begin{column}{0.55\textwidth}
            \textbf{Contexto e Requisitos:}
            \begin{itemize}
                \item Registro diário de gastos do usuário.
                \item Geração de relatórios mensais de consumo por categorias.
                \item Vínculo de membros familiares (dados isolados por grupo).
                \item Inteligência Artificial para apoiar no planejamento orçamentário mensal.
            \end{itemize}
        \end{column}
        \begin{column}{0.45\textwidth}
            \textbf{Desafio Metodológico:}
            Como uma equipe de 5 alunos pode estimar com precisão o esforço inicial, definir a Sprint 1 e planejar prazos de forma robusta e baseada em dados estatísticos?
        \end{column}
    \end{columns}
\end{frame}

\section{Método Teórico Clássico}
\begin{frame}{Método Teórico: Estimativa Relativa e Fibonacci}
    \begin{itemize}
        \item \textbf{Story Points (Pontos de História):} Medida abstrata que quantifica esforço, complexidade e incerteza técnica de tarefas.
        \item \textbf{Sequência de Fibonacci ($1, 2, 3, 5, 8, 13$):} Reduz a falsa precisão de estimativas lineares. Quanto maior a tarefa, maior a incerteza.
        \item \textbf{Planning Poker:} Reunião de consenso onde cada desenvolvedor apresenta seu palpite sem sofrer influência inicial do grupo.
    \end{itemize}
    \begin{block}{Parâmetros Estatísticos de Incerteza}
        Calculamos a Média Amostral ($\bar{X}$), o Desvio Padrão Amostral ($s$) e o Erro Padrão ($SE$) dos palpites para medir a dispersão e incerteza da equipe.
    \end{block}
\end{frame}

\begin{frame}{Exemplo Teórico Analítico (Passo a Passo)}
    \textbf{Tarefa:} "Integrar API do Gemini para gerar relatórios preditivos"
    \begin{itemize}
        \item Estimativas fornecidas por 5 desenvolvedores: $5, 8, 5, 8, 8$ Story Points.
    \end{itemize}
    
    \textbf{Passo 1: Média Amostral ($\bar{X}$):}
    $$\bar{X} = \frac{5 + 8 + 5 + 8 + 8}{5} = \frac{34}{5} = 6.8 \text{ SP}$$
    
    \textbf{Passo 2: Desvio Padrão Amostral ($s$):}
    $$s = \sqrt{\frac{(5-6.8)^2 + (8-6.8)^2 + (5-6.8)^2 + (8-6.8)^2 + (8-6.8)^2}{5 - 1}}$$
    $$s = \sqrt{\frac{3.24 + 1.44 + 3.24 + 1.44 + 1.44}{4}} = \sqrt{2.7} \approx 1.64 \text{ SP}$$
    
    \textbf{Passo 3: Erro Padrão ($SE$):}
    $$SE = \frac{s}{\sqrt{N}} = \frac{1.64}{\sqrt{5}} \approx 0.73 \text{ SP}$$
\end{frame}

\section{Abordagem Computacional}
\begin{frame}[fragile]{Demonstração Computacional Equivalente}
    O código Python a seguir realiza o mesmo cálculo de incerteza feito analiticamente:
\begin{lstlisting}
import numpy as np

# Estimativas dos 5 alunos
estimativas = np.array([5, 8, 5, 8, 8])
n = len(estimativas)

# Calculos Estatisticos
media = np.mean(estimativas)
desvio_padrao = np.std(estimativas, ddof=1) # ddof=1 garante desvio padrao amostral (N-1)
erro_padrao = desvio_padrao / np.sqrt(n)

print(f"Media Amostral: {media:.2f} SP")
print(f"Desvio Padrao Amostral: {desvio_padrao:.2f} SP")
print(f"Erro Padrao: {erro_padrao:.2f} SP")
\end{lstlisting}
    \begin{block}{Conexão Computacional}
        A saída do código valida as equações de Média Amostral ($6.80$ SP), Desvio Padrão ($1.64$ SP) e Erro Padrão ($0.73$ SP).
    \end{block}
\end{frame}

\begin{frame}{Simulação Monte Carlo para Prazos}
    \begin{itemize}
        \item O backlog total do MVP foi fechado em \textbf{80 Story Points (SP)}.
        \item A velocidade da equipe por sprint segue uma distribuição normal com média de \textbf{15 SP} e desvio padrão de \textbf{3.5 SP}.
        \item Realizamos $10.000$ simulações em Python para prever o prazo de entrega.
    \end{itemize}
    \begin{block}{Resultados obtidos via simulação}
        \begin{itemize}
            \item \textbf{50\% de Confiança (Mais Provável):} Conclusão em até 5.3 Sprints.
            \item \textbf{80\% de Confiança (Recomendado):} Conclusão em até 5.8 Sprints.
            \item \textbf{90\% de Confiança (Pessimista):} Conclusão em até 6.1 Sprints.
        \end{itemize}
    \end{block}
\end{frame}

\section{Planejamento e Execução}
\begin{frame}{Trilha ADS: Divisão das Tarefas da Sprint 1}
    Para a primeira sprint (2 semanas), a equipe de ADS planejou \textbf{16 SP} de capacidade:
    \begin{table}[]
        \tiny
        \centering
        \begin{tabular}{lllc}
            \toprule
            \textbf{Responsável} & \textbf{Papel Principal} & \textbf{Atividades da Sprint 1} & \textbf{Esforço (SP)} \\
            \midrule
            Aluno A & Infra & Setup Next.js 16 + Docker + Prisma PostgreSQL & 3 SP \\
            Aluno B & Segurança & Autenticação Next-Auth + Lógica de Vínculo Familiar & 5 SP \\
            Aluno C & UI/UX & Tela móvel de cadastro e Histórico de Gastos & 3 SP \\
            Aluno D & Backend & API CRUD de Categorias de gastos e deleção lógica & 2 SP \\
            Aluno E & IA Integrador & Conexão SDK Gemini + Loader Assíncrono (Skeleton) & 3 SP \\
            \bottomrule
        \end{tabular}
        \caption{Divisão equilibrada da carga de trabalho para 5 alunos de ADS}
    \end{table}
\end{frame}

\begin{frame}{Trilha Técnica: Divisão Alternativa de Tarefas}
    Para o curso Técnico em Informática para Internet, o foco é frontend e consumo de APIs (13 SP):
    \begin{table}[]
        \tiny
        \centering
        \begin{tabular}{lllc}
            \toprule
            \textbf{Responsável} & \textbf{Papel Principal} & \textbf{Atividades da Sprint 1} & \textbf{Esforço (SP)} \\
            \midrule
            Aluno 1 & Web Design & Prototipagem e estilização responsiva do painel (Tailwind) & 3 SP \\
            Aluno 2 & Integrator & Roteamento dinâmico (App Router) e proteção simples de rotas & 3 SP \\
            Aluno 3 & JavaScript & Formulário de gastos, validações client-side e lista local & 2 SP \\
            Aluno 4 & Visualização & Gráfico de pizza de gastos mensais por categorias (Recharts) & 2 SP \\
            Aluno 5 & APIs / IA & Envio de requisições HTTP e exibição das dicas de IA na tela & 3 SP \\
            \bottomrule
        \end{tabular}
        \caption{Divisão de tarefas de nível técnico ajustada para 5 alunos}
    \end{table}
\end{frame}

\begin{frame}{Arquitetura do FamilyFinance.AI}
    Desenvolvimento baseado em \textbf{Next.js 16} (React 19) e banco de dados isolado por família:
    \begin{columns}
        \begin{column}{0.5\textwidth}
            \textbf{Camada Web (Next.js 16):}
            \begin{itemize}
                \item App Router para organização.
                \item Server Components (RSC) para leitura.
                \item Server Actions para mutações.
                \item Vercel AI SDK para integração.
            \end{itemize}
        \end{column}
        \begin{column}{0.5\textwidth}
            \textbf{Isolamento Multi-tenant:}
            \begin{itemize}
                \item Toda transação e categoria possui o vínculo obrigatório a um \texttt{familyId}.
                \item Consultas do Prisma são filtradas pela sessão do usuário logado no servidor.
            \end{itemize}
        \end{column}
    \end{columns}
\end{frame}

\section{Metas e Exercícios}
\begin{frame}{Guia de OKR Individual e Gestão de Metas}
    Cada aluno utilizará o modelo de OKR e fará o auto-monitoramento diário:
    \begin{itemize}
        \item \textbf{Daily Check-in:} O que fiz ontem? O que farei hoje? Algum impedimento no meu progresso?
        \item \textbf{Foco na Meta da Sprint:} Garantir que as entregas de infra e fluxo básico de autenticação/CRUD funcionem sem atritos.
    \end{itemize}
    \begin{alertblock}{Impedimentos Coletivos}
        Qualquer gargalo técnico (ex: conflitos de tipagem ou portas do Docker) deve ser resolvido imediatamente em reuniões rápidas e cooperativas.
    \end{alertblock}
\end{frame}

\begin{frame}{Exercícios Propostos para os Alunos}
    \begin{enumerate}
        \item \textbf{Exercício Analítico:} Uma tarefa recebeu as estimativas $3, 5, 3, 5, 8$ SP de 5 alunos. Calcule a Média Amostral ($\bar{X}$), o Desvio Padrão Amostral ($s$) e o Erro Padrão ($SE$).
        \item \textbf{Exercício de Decisão:} Com base no gráfico ECDF, se a diretoria exigir 80\% de confiança na data de entrega do MVP, qual deve ser o prazo acordado (em número de sprints)?
        \item \textbf{Desafio de Código:} Altere as propriedades do script Python para uma velocidade média pessimista de 12 SP e desvio padrão de 4.5 SP. Plote a nova ECDF e aponte o novo prazo para 80\% de confiança.
    \end{enumerate}
\end{frame}

\end{document}
